\documentclass[11pt,a4paper]{article}
\usepackage{amsmath,amssymb,amsthm}
\usepackage{geometry}
\geometry{margin=1in}
\usepackage{hyperref}
\usepackage{enumitem}

\newtheorem{theorem}{Theorem}[section]
\newtheorem{lemma}[theorem]{Lemma}
\newtheorem{proposition}[theorem]{Proposition}
\newtheorem{corollary}[theorem]{Corollary}
\newtheorem{definition}[theorem]{Definition}
\theoremstyle{remark}
\newtheorem{remark}[theorem]{Remark}
\newtheorem{problem}{Problem}[section]

\newcommand{\cL}{\mathcal{L}}
\newcommand{\cH}{\mathcal{H}}
\newcommand{\cB}{\mathcal{B}}
\newcommand{\cC}{\mathcal{C}}
\newcommand{\bbC}{\mathbb{C}}
\newcommand{\bbN}{\mathbb{N}}
\newcommand{\bbR}{\mathbb{R}}
\newcommand{\bbZ}{\mathbb{Z}}
\DeclareMathOperator{\spec}{spec}
\DeclareMathOperator{\tr}{tr}
\DeclareMathOperator{\detn}{det}

\title{\bfseries Spectral Analysis of the Limiting Hilbert–Pólya Operator\\
\large A Research Programme}
\author{}
\date{}

\begin{document}
\maketitle

\begin{abstract}
The harmonic–categorical approach to the Riemann zeta function has yielded a sequence of finite self‑adjoint matrices $H_k$ whose spectral measures converge to the zero‑counting measure of the critical zeros. The programme outlined here addresses the rigorous construction and analysis of the limiting self‑adjoint operator $H$ obtained from this sequence. The central tasks are to identify the domain and action of $H$, to prove that its spectrum consists precisely of the ordinates of the non‑trivial zeros (with multiplicities), and to establish the precise connection between $H$ and the transfer operator families $\cL_{1/2+it}$. A successful completion would constitute a proof of the Riemann Hypothesis.
\end{abstract}

\section{Background}
We summarise the essential results obtained so far (see \cite{findings} for details).

\subsection{The Greedy Transfer Operator}
From the greedy harmonic tree one constructs a transfer operator $\cL_s$ acting on the Hardy space $\cH = H^2(\mathbb{D})$, trace‑class for $\operatorname{Re}s>\tfrac12$, whose Fredholm determinant satisfies
\begin{equation}\label{eq:det}
\det\nolimits_{(1)}(I-\cL_s)=\frac{\zeta(s)}{\zeta(2s)}\,e^{P(s)},
\end{equation}
with $P(s)$ entire and non‑vanishing on the critical strip. Hence the non‑trivial zeros of $\zeta(s)$ are exactly the parameters where $1$ is an eigenvalue of $\cL_s$.

\subsection{Finite‑Dimensional Approximations}
For each integer $k\ge1$ we construct a finite‑dimensional truncation $\cL_s^{(k)}$ of $\cL_s$, acting on a subspace $\cH_k\subset\cH$ of dimension $d_k\to\infty$. The determinants converge uniformly on compact sets:
\[
\det(I-\cL_s^{(k)})\longrightarrow\det(I-\cL_s).
\]
On the critical line we set $A_k(t)=\cL_{1/2+it}^{(k)}$. By polar decomposition $A_k(t)=U_k(t)|A_k(t)|$ we obtain a family of unitary matrices $U_k(t)$, analytic in $t$ on intervals free of singularities. The finite‑dimensional Hamiltonians are defined as
\[
H_k = i\,\frac{d}{dt}\log U_k(t)\Big|_{t=0},
\]
where the branch of the logarithm is chosen so that $H_k$ is Hermitian. The eigenvalues of $H_k$ are precisely the numbers $\gamma$ for which $\det(I-A_k(\gamma))=0$, i.e.\ approximations to the ordinates of the critical zeros.

\subsection{Convergence of Spectral Measures}
Let $\mu_k = \frac{1}{d_k}\sum_{\lambda\in\spec(H_k)}\delta_\lambda$. Then $\mu_k$ converges vaguely to the measure $\mu = \sum_{\gamma} \delta_\gamma$, where the sum is over all ordinates of non‑trivial zeros on the critical line (counted with multiplicity). This follows from the convergence of determinants and the standard link between zeros and eigenvalues of $H_k$.

\section{The Limiting Operator $H$}
The sequence $\{H_k\}$ does not converge in a strong operator sense because their dimensions increase. However, the convergence of spectral measures suggests that there exists a self‑adjoint operator $H$ on a Hilbert space, obtained as a limit of the $H_k$ in a resolvent sense, such that $\spec(H)=\operatorname{supp}\mu$. The central goals of the programme are:
\begin{enumerate}[label=(\roman*)]
    \item Give a constructive definition of $H$ acting on a concrete Hilbert space (e.g.\ a space of functions on the boundary of the greedy tree, or an $L^2$‑space over the critical line).
    \item Prove that $H$ is self‑adjoint with the prescribed spectrum, including the correct multiplicities.
    \item Relate $H$ directly to the infinitesimal generator of the unitary group associated with the transfer operators $\cL_{1/2+it}$, thereby closing the circle with Problem~9 of the original programme.
\end{enumerate}

\section{Open Problems}
The following problems are designed to achieve the above goals. They are concrete, well‑posed, and build on the existing analytic framework.

\begin{problem}[Direct limit of truncated operators]
Construct a Hilbert space $\cH_\infty$ as an inductive limit of the subspaces $\cH_k$, and define a self‑adjoint operator $H$ on $\cH_\infty$ such that the finite‑dimensional projections of $H$ reproduce the $H_k$ in the sense of spectral measures. Prove that the resolvent $(H-z)^{-1}$ is compact.
\end{problem}

\begin{problem}[Generator from the unitary group]
Recall that the problem of defining $U_t = \cL_{1/2+it}\,\cL_{1/2-it}^{-1}$ was obstructed at the zeros. However, on the orthogonal complement of the eigenspaces where $1$ is an eigenvalue, $U_t$ exists. Construct a maximal sectorially defined unitary group and use it to obtain a self‑adjoint generator $H$ via a suitable boundary‑value construction. Prove that the complementary spectrum of $H$ corresponds exactly to the Riemann zeros.
\end{problem}

\begin{problem}[Alternative: Dirac operator on the tree]
Using the spin structure and Laplacian on the greedy tree, define a Dirac operator $D$ and a twisted version $D_s$ that encodes the transfer operator weights. Show that the spectrum of the square $D_{1/2}^2$ (or a related self‑adjoint operator) is purely discrete and equals $\{\gamma^2:\zeta(1/2+i\gamma)=0\}$. This would provide a geometric realisation akin to the Connes–Consani programme.
\end{problem}

\begin{problem}[Spectral flow and holonomy]
Study the spectral flow of the family $\cL_{1/2+it}$ along the critical line. Define the holonomy unitary on the determinant bundle and derive a self‑adjoint operator from its logarithm. Prove that the holonomy eigenvalues are exactly the zero ordinates, with multiplicities given by the winding numbers of the eigenvalues of $\cL_{1/2+it}$.
\end{problem}

\begin{problem}[Rigorous proof that the limit spectrum is exactly the critical zeros]
Assuming the Riemann Hypothesis, the support of $\mu$ is known. But even without RH, prove that every accumulation point of the eigenvalues of $H_k$ is either an ordinate of a critical zero or a limit of such ordinates. Show that no additional real numbers can appear in the limit spectrum. This would establish that the limiting $H$ captures exactly the critical line zeros, with no extraneous spectrum.
\end{problem}

\begin{problem}[Connection with the explicit formula]
Derive an explicit formula for the trace of $(H-z)^{-1}$ in terms of the logarithmic derivative of $\zeta(s)$ and the entire factor $P(s)$. Use this to prove that the spectral measure of $H$ is exactly the zero‑counting measure, without assuming RH. This would be a spectral interpretation of the Guinand–Weil formula within the harmonic sine framework.
\end{problem}

\begin{problem}[Numerical verification and refinements]
Implement a large‑scale numerical computation of $H_k$ for increasing $k$ and verify the convergence of the spectral measures. Investigate the statistics of eigenvalue spacings and compare with predictions from random matrix theory. This would bolster confidence in the construction and may reveal finer properties of $H$.
\end{problem}

\section{Expected Outcomes}
A complete solution to these problems would yield:
\begin{itemize}
\item A self‑adjoint operator $H$, defined explicitly in terms of the greedy harmonic tree or modular tensor categories, with $\spec(H)=\{\gamma:\zeta(1/2+i\gamma)=0\}$.
\item A proof that the Riemann Hypothesis is equivalent to the absence of any other spectrum (which would be a consequence if the construction is tight).
\item A new spectral proof of the functional equation for $\zeta(s)$ via the properties of $H$.
\end{itemize}
Even partial progress will radically deepen our understanding of the spectral nature of the zeros and the role of the harmonic sine transform in number theory.

\begin{thebibliography}{9}
\bibitem{programme} V.~Geere et al., \emph{A Harmonic–Categorical Approach to the Riemann Zeta Function: Research Programme}, 2026.
\bibitem{findings} ---, \emph{A Harmonic–Categorical Operator for the Riemann Zeta Function: Findings of the Research Programme}, 2026.
\bibitem{Geere2026a} V.~Geere, \emph{A Purely Geometric Reconstruction of the Sine Function}, \url{https://victorgeere.co.za/maths/sine-harmonic.html}.
\bibitem{Geere2026b} V.~Geere, \emph{On the Correlation Structure of a Greedy Harmonic Decomposition of the Sine Function}, \url{https://victorgeere.co.za/maths/greedy-harmonic-decomposition.html}.
\bibitem{Geere2026c} V.~Geere, \emph{The Quantum‑Egyptian Functional Equation}, v8, \url{https://victorgeere.co.za/maths/quantum-egyptian-functional-equation-v8.html}.
\bibitem{Mayer} D.~H.~Mayer, \emph{The thermodynamic formalism approach to Selberg's zeta function for $\mathrm{PSL}(2,\mathbb{Z})$}, Bull. Amer. Math. Soc. (N.S.) 25 (1991), 55--60.
\bibitem{Efrat} I.~Efrat, \emph{Dynamics of the continued fraction map and the spectral theory of $\mathrm{SL}(2,\mathbb{Z})$}, Invent. Math. 114 (1993), 207--218.
\bibitem{ConnesConsani} A.~Connes, C.~Consani, \emph{Schemes over $\mathbb{F}_1$ and zeta functions}, Compos. Math. 146 (2010), 1383--1415.
\end{thebibliography}

\end{document}